% GENERATED FILE. Edit canonical lessons, then run npm run generate:books.
% canonical-source-hash: fnv1a64:d3e2e3d2265797eb
% canonical-lessons: BN-C354-gobeshona, BN-C354-porikkha, BN-C354-tottvo, BN-C354-abishkar-kora, BN-C354-onusondhan-kora

\chapter{Finding Out}
\label{ch:bn-finding-out}
\hlchaptermodality{354}

\begin{chapteropening}
\textbf{By the end of this chapter:} \emph{I can say what a study looked for and what it found.}

Complete the last lesson of chapter 354: I can say what a study looked for and what it found.
\end{chapteropening}

\section[\texorpdfstring{\bn{গবেষণা} (gôbeshônā)}{গবেষণা (gôbeshônā)}]{\bn{গবেষণা} (gôbeshônā) — research}

\label{lesson:BN-C354-gobeshona}



\begin{quote}
Before the new one: say the Bengali for in common use, then the Bengali for to obey.
\end{quote}

\subsection*{You'll want to know: \bn{গবেষণা}}
\textbf{\bn{গবেষণা}} — \emph{gôbeshônā} — \textquotedblleft{}research\textquotedblright{}.

\begin{quote}
\bn{মীরা} \bn{বিশ্ববিদ্যালয়ে} \bn{গবেষণা} \bn{করে।}
\end{quote}

\begin{quote}
\emph{mīrā bishvôbidyālôye gôbeshônā kôre.}
\end{quote}

\begin{quote}
Mira does research at the university.
\end{quote}

\begin{cousinweb}
\textbf{\bn{গবেষণা}} (\emph{gôbeshônā}) is Sanskrit \emph{gaveṣaṇā}, beside \emph{gaveṣaṇa}, seeking or searching (Monier-Williams), which files the verb \emph{gaveṣ}, to seek, among the compounds of \emph{go}, a cow. Bengali uses it for scholarly research.
\end{cousinweb}

\begin{grammarlens}[title={What you've built}]
One more word for giving reasons and opinions.
\end{grammarlens}

\subsection*{Your turn}
\begin{itemize}
\raggedright
  \item \emph{Say it:} \emph{gôbeshônā}
  \item \emph{Say it:} \emph{gôbeshônā}, once more
  \item \emph{Say it:} \emph{mīrā bishvôbidyālôye gôbeshônā kôre.} — Mira does research at the university.
  \item \emph{Recall:} say the Bengali for in common use, then the Bengali for to obey, then say \emph{gôbeshônā} again
\end{itemize}

\begin{tcolorbox}[breakable,colback=teal!4,colframe=teal!35!black,title={Before you move on}]
What does \bn{গবেষণা} mean? (\textquotedblleft{}Research\textquotedblright{}.) Say it once more.
\end{tcolorbox}
\section[\texorpdfstring{\bn{পরীক্ষা} (pôrīkkhā)}{পরীক্ষা (pôrīkkhā)}]{\bn{পরীক্ষা} (pôrīkkhā) — a test; an experiment; an exam}

\label{lesson:BN-C354-porikkha}



\begin{quote}
Before the new one: say the Bengali for to obey, then the Bengali for research.
\end{quote}

\subsection*{You'll want to know: \bn{পরীক্ষা}}
\textbf{\bn{পরীক্ষা}} — \emph{pôrīkkhā} — \textquotedblleft{}a test; an experiment; an exam\textquotedblright{}.

\begin{quote}
\bn{বিজ্ঞানীরা} \bn{একটা} \bn{নতুন} \bn{পরীক্ষা} \bn{করলেন।}
\end{quote}

\begin{quote}
\emph{bigyānīrā ekṭā nôtun pôrīkkhā kôrlen.}
\end{quote}

\begin{quote}
The scientists carried out a new experiment.
\end{quote}

\begin{cousinweb}
\textbf{\bn{পরীক্ষা}} (\emph{pôrīkkhā}) is Sanskrit \emph{parīkṣā}, looking all round, inspection or test (Monier-Williams). The same word serves for a school exam and a laboratory experiment.
\end{cousinweb}

\begin{grammarlens}[title={What you've built}]
One more word for giving reasons and opinions.
\end{grammarlens}

\subsection*{Your turn}
\begin{itemize}
\raggedright
  \item \emph{Say it:} \emph{pôrīkkhā}
  \item \emph{Say it:} \emph{pôrīkkhā}, once more
  \item \emph{Say it:} \emph{bigyānīrā ekṭā nôtun pôrīkkhā kôrlen.} — The scientists carried out a new experiment.
  \item \emph{Recall:} say the Bengali for to obey, then the Bengali for research, then say \emph{pôrīkkhā} again
\end{itemize}

\begin{tcolorbox}[breakable,colback=teal!4,colframe=teal!35!black,title={Before you move on}]
What does \bn{পরীক্ষা} mean? (\textquotedblleft{}A test; an experiment; an exam\textquotedblright{}.) Say it once more.
\end{tcolorbox}
\section[\texorpdfstring{\bn{তত্ত্ব} (tôttvô)}{তত্ত্ব (tôttvô)}]{\bn{তত্ত্ব} (tôttvô) — a theory}

\label{lesson:BN-C354-tottvo}



\begin{quote}
Before the new one: say the Bengali for research, then the Bengali for an experiment.
\end{quote}

\subsection*{You'll want to know: \bn{তত্ত্ব}}
\textbf{\bn{তত্ত্ব}} — \emph{tôttvô} — \textquotedblleft{}a theory\textquotedblright{}.

\begin{quote}
\bn{এই} \bn{তত্ত্বের} \bn{কোনো} \bn{প্রমাণ} \bn{নেই।}
\end{quote}

\begin{quote}
\emph{ei tôttver kono prômān nei.}
\end{quote}

\begin{quote}
There is no proof for this theory.
\end{quote}

\begin{cousinweb}
\textbf{\bn{তত্ত্ব}} (\emph{tôttvô}) is Sanskrit \emph{tattva}, 'thatness', the true state of things (Monier-Williams). Modern Bengali uses it for a theory, an account of how things really are.
\end{cousinweb}

\begin{grammarlens}[title={What you've built}]
One more word for giving reasons and opinions.
\end{grammarlens}

\subsection*{Your turn}
\begin{itemize}
\raggedright
  \item \emph{Say it:} \emph{tôttvô}
  \item \emph{Say it:} \emph{tôttvô}, once more
  \item \emph{Say it:} \emph{ei tôttver kono prômān nei.} — There is no proof for this theory.
  \item \emph{Recall:} say the Bengali for research, then the Bengali for an experiment, then say \emph{tôttvô} again
\end{itemize}

\begin{tcolorbox}[breakable,colback=teal!4,colframe=teal!35!black,title={Before you move on}]
What does \bn{তত্ত্ব} mean? (\textquotedblleft{}A theory\textquotedblright{}.) Say it once more.
\end{tcolorbox}
\section[\texorpdfstring{\bn{আবিষ্কার} \bn{করা} (ābishkār kôrā)}{আবিষ্কার করা (ābishkār kôrā)}]{\bn{আবিষ্কার} \bn{করা} (ābishkār kôrā) — to discover; to invent}

\label{lesson:BN-C354-abishkar-kora}



\begin{quote}
Before the new one: say the Bengali for an experiment, then the Bengali for a theory.
\end{quote}

\subsection*{You'll want to know: \bn{আবিষ্কার} \bn{করা}}
\textbf{\bn{আবিষ্কার} \bn{করা}} — \emph{ābishkār kôrā} — \textquotedblleft{}to discover; to invent\textquotedblright{}.

\begin{quote}
\bn{বিজ্ঞানীরা} \bn{নদীতে} \bn{একটা} \bn{নতুন} \bn{মাছ} \bn{আবিষ্কার} \bn{করলেন।}
\end{quote}

\begin{quote}
\emph{bigyānīrā nôdīte ekṭā nôtun māchh ābishkār kôrlen.}
\end{quote}

\begin{quote}
The scientists discovered a new fish in the river.
\end{quote}

Past (he/she …): \bn{তিনি} \bn{একটা} \bn{নতুন} \bn{পদ্ধতি} \bn{আবিষ্কার} \bn{করলেন}, \emph{tini ekṭā nôtun pôddhôti ābishkār kôrlen}.

\begin{cousinweb}
\textbf{\bn{আবিষ্কার} \bn{করা}} (\emph{ābishkār kôrā}) is Sanskrit \emph{āviṣkāra}, from \emph{āviṣ} 'openly, in sight' and \emph{kāra} 'making' (Monier-Williams), with \bn{করা}: to bring into view what was hidden, or to make a thing that did not exist.
\end{cousinweb}

\begin{grammarlens}[title={What you've built}]
One more word for giving reasons and opinions.
\end{grammarlens}

\subsection*{Your turn}
\begin{itemize}
\raggedright
  \item \emph{Say it:} \emph{ābishkār kôrā}
  \item \emph{Say it:} \emph{ābishkār kôrā}, once more
  \item \emph{Say it:} \emph{bigyānīrā nôdīte ekṭā nôtun māchh ābishkār kôrlen.} — The scientists discovered a new fish in the river.
  \item \emph{Recall:} say the Bengali for an experiment, then the Bengali for a theory, then say \emph{ābishkār kôrā} again
\end{itemize}

\begin{tcolorbox}[breakable,colback=teal!4,colframe=teal!35!black,title={Before you move on}]
What does \bn{আবিষ্কার} \bn{করা} mean? (\textquotedblleft{}To discover; to invent\textquotedblright{}.) Say it once more.
\end{tcolorbox}
\section[\texorpdfstring{\bn{অনুসন্ধান} \bn{করা} (ônusôndhān kôrā)}{অনুসন্ধান করা (ônusôndhān kôrā)}]{\bn{অনুসন্ধান} \bn{করা} (ônusôndhān kôrā) — to investigate; to look into}

\label{lesson:BN-C354-onusondhan-kora}



\begin{quote}
Before the new one: say the Bengali for a theory, then the Bengali for to discover.
\end{quote}

\subsection*{You'll want to know: \bn{অনুসন্ধান} \bn{করা}}
\textbf{\bn{অনুসন্ধান} \bn{করা}} — \emph{ônusôndhān kôrā} — \textquotedblleft{}to investigate; to look into\textquotedblright{}.

\begin{quote}
\bn{আমরা} \bn{দুর্ঘটনার} \bn{কারণ} \bn{অনুসন্ধান} \bn{করব।}
\end{quote}

\begin{quote}
\emph{āmrā durghôṭônār kārôn ônusôndhān kôrbô.}
\end{quote}

\begin{quote}
We will look into the cause of the accident.
\end{quote}

Past (he/she …): \bn{তিনি} \bn{কারণটা} \bn{অনুসন্ধান} \bn{করলেন}, \emph{tini kārônṭā ônusôndhān kôrlen}.

\begin{cousinweb}
\textbf{\bn{অনুসন্ধান} \bn{করা}} (\emph{ônusôndhān kôrā}) is Sanskrit \emph{anusaṃdhāna}, investigation or close inquiry (Monier-Williams), with \bn{করা}.
\end{cousinweb}

\begin{grammarlens}[title={What you've built}]
One more word for giving reasons and opinions.
\end{grammarlens}

\subsection*{Your turn}
\begin{itemize}
\raggedright
  \item \emph{Say it:} \emph{ônusôndhān kôrā}
  \item \emph{Say it:} \emph{ônusôndhān kôrā}, once more
  \item \emph{Say it:} \emph{āmrā durghôṭônār kārôn ônusôndhān kôrbô.} — We will look into the cause of the accident.
  \item \emph{Recall:} say the Bengali for a theory, then the Bengali for to discover, then say \emph{ônusôndhān kôrā} again
\end{itemize}

\begin{tcolorbox}[breakable,colback=teal!4,colframe=teal!35!black,title={Before you move on}]
What does \bn{অনুসন্ধান} \bn{করা} mean? (\textquotedblleft{}To investigate; to look into\textquotedblright{}.) Say it once more.
\end{tcolorbox}
